\chapter{Introdução}

O termo ‘câncer’ tem origem na palavra grega Karkinos (caranguejo), pois o tumor característico da doença deforma a pele sobre os vasos sanguíneos, remetendo a figura de um caranguejo.
Egípcios e gregos fizeram os primeiros registros sobre tumores nos seios, tratando a doença com amputações e remédios que incluíam miolos de vaca e excremento de vespa.
Acreditava-se também que o sangue menstrual era capaz de subir às mamas e transformar-se em leite, assim como causar tumores ao encaroçar-se nos seios.\cite{inca2018mulher}

O câncer de mama é a multiplicação desordenada de células anormais nas mamas, resultando em um tumor maligno que pode invadir outros tecidos e órgãos, e atualmente é a neoplasia maligna mais incidente entre as mulheres em todo o mundo e a primeira causa de morte por câncer entre mulheres no Brasil, com patamares diferenciados entre as regiões. \cite{ms2026cancermama} Sua alta incidência e prevalência tornaram-no um dos principais desafios de saúde pública no país, destacando a necessidade de estratégias eficazes para sua detecção precoce e manejo \cite{Pinheiro2024}.

Em 2025, foram estimados 73.610 casos novos e o estado de Minas Gerais (MG) concentra 7.670 desses casos, o que corresponde a 10,42\% do total nacional. No que se refere à mortalidade, dos 20.165 óbitos registrados, 1.829 ocorreram em MG, representando 9.07\% do total \cite{Inca2024}. Esses números revelam que o estado tem uma participação expressiva tanto na incidência quanto nos óbitos, o que levanta questionamentos importantes sobre os fatores locais que podem estar contribuindo para esses índices elevados. 

Os problemas relacionados ao câncer de mama vão além da doença em si e incluem uma extensa rede de contribuintes. Entre os fatores comportamentais e metabólicos, destacam-se o consumo de álcool, a obesidade, padrões reprodutivos e hormonais, além da alimentação e o tabagismo, possivelmente \cite{Alves2022}. No campo social e econômico, estudos apontam que diagnósticos em estágio avançado podem estar vinculados à baixa escolaridade, à raça/cor não branca e à menor disponibilidade de equipamentos e serviços de mamografia \cite{Dantas2021}. Outros aspectos relevantes incluem a cobertura da Atenção Primária à Saúde (APS) e as desigualdades territoriais, bem como as interseções entre gênero e renda, também pesam fatores de ambiente e ocupação. Do ponto de vista do sistema de saúde, aspectos centrais incluem a cobertura e periodicidade do rastreamento por mamografia, o tempo entre sintomas, diagnóstico e início do tratamento, o estágio em que a doença é detectada, a disponibilidade e distância de serviços de oncologia (exames, mamografias, cirurgia, radioterapia), falta de conhecimento sobre sinais de alerta e importância do autoexame. Nesse sentido,  pesquisas evidenciam a persistência de barreiras internas e externas à implementação do rastreamento mamográfico no Sistema Único de Saúde (SUS) \cite{Sala2025}. 

Mesmo com esses avanços analíticos e com o progresso das políticas públicas de rastreamento e conscientização como o Outubro Rosa e a ampliação da cobertura de rastreamento no SUS, ainda persistem barreiras no acesso à informação e à prevenção da doença, visto que, essas movimentações não estão sendo suficientes para diminuir os fatores que tornam câncer de mama como uma das doenças que mais causa a morte de mulheres no país. Apesar da relevância das múltiplas variáveis relacionadas à incidência do câncer de mama em MG, existem lacunas de estudos que integrem diferentes fatores de forma correlacionada. Essa barreira dificulta a identificação de componentes que possam explicar as variações regionais da doença e limita a compreensão do fenômeno em sua totalidade.

Ao delimitar o campo de estudo para uma região específica, como o estado de MG e seus 853 municípios, torna-se possível realizar o mapeamento espacial e temporal dos casos. Esse olhar integrado permitirá compreender não apenas os números de incidência e mortalidade, mas também alguns dos possíveis obstáculos enfrentados pelas pessoas no cuidado com a própria saúde. Dessa forma, será possível identificar como determinadas características afetam de maneira diferenciada as regiões do estado, quais fatores locais podem estar ou não mais diretamente relacionados à incidência da doença e onde é necessário mais estratégias de prevenção e cuidado.

% frequentemente substituída por uma expectativa de melhora espontânea que pode não se concretizar.

Diante deste cenário, este trabalho busca responder quais fatores socioeconômicos municipais, condições de acesso à infraestrutura de saúde e características epidemiológicas das pacientes se associam à vulnerabilidade ao câncer de mama em MG. Assume-se que essa vulnerabilidade não é distribuída de forma homogênea, estando fortemente correlacionada à baixa oferta de serviços de média e alta complexidade nas macrorregiões de saúde, bem como a indicadores de baixa escolaridade e renda, ano de diagnósticos, hospitais disponíveis que retardam o diagnóstico precoce.

\section{Objetivo Geral}

Analisar e identificar padrões e variações referente a incidência do câncer de mama no Brasil, com foco no estado de MG e em seus municípios, relacionando variáveis de caráter socioeconômico e demográfico, de modo a subsidiar a pesquisa e a compreensão dos fatores associados à doença.

Esse objetivo principal será buscado por meio dos seguintes objetivos específicos:

\begin{enumerate}
\item "Realizar o pré-processamento de dados de múltiplas bases públicas, abrangendo as etapas de seleção, limpeza, integração e transformação para consolidar um conjunto de informações unificado
\item Construir uma análise multivariada sobre os casos de câncer de mama em MG, identificando a posição do estado  em relação aos demais e correlacionando o número de casos, mortalidade, acesso ao tratamento, ocupação populacional e condições sociais dos municípios.
\item Utilizar técnicas computacionais ao longo do processo, como mineração de dados, algoritmos de agrupamento e métodos de visualização gráfica (gráficos e mapas de calor) e teoria dos grafos para auxiliar na identificação de padrões e relações entre variáveis.
\item Realizar um mapeamento descritivo capaz de evidenciar lacunas, desigualdades e fatores que influenciam ou não o aumento ou a diminuição dos casos, para assim, oferecer subsídios analíticos para estudos futuros e para o planejamento em saúde e de políticas públicas no Brasil.
\end{enumerate}

\chapter{Estrutura da Monografia}

Esse texto foi estruturado em 5 capítulos. O primeiro capítulo é a introdução, onde são apresentados o tema, a justificativa, os objetivos e a estrutura do trabalho. O segundo capítulo é o referencial teórico, onde são revisados os principais conceitos e teorias da literatura no contexto deste trabalho. O terceiro capítulo é a metodologia, onde são descritos os métodos utilizados para a coleta, tratamento e análise dos dados. O quarto capítulo é dedicado aos resultados, onde são apresentados os principais achados da pesquisa. O quinto capítulo é a conclusão, onde há as dicussões e apresentadas as principais considerações finais do trabalho.

% \section{Figuras e Tabelas}
% Nos elementos flutuantes, as legendas devem estar alinhadas à esquerda com o comando \textbf{minipage}.

% \begin{verbatim}
% \begin{table}[!ht]
%     \centering
%     \begin{minipage}{0.7\textwidth}     <<<<<<<
%         \caption{\label{tabela:ComparativoFrameworks}
%         Comparação dos Frameworks}
%         \resizebox{\textwidth}{!}{
%             [...]
%         }
%         \caption*{\footnotesize Fonte: Elaborado pelo autor, 2023.}
%     \end{minipage}
% \end{table}
% \end{verbatim}

% \begin{table}[!ht]
%     \centering
%     \begin{minipage}{0.7\textwidth}
%         \caption{\label{tabela:ComparativoFrameworks}
%             Comparação dos Frameworks}
%         \resizebox{\textwidth}{!}{
%             \begin{tabular}{|l|c|c|c|}
%                 \hline
%                                        & \textbf{MapReduce} & \textbf{Spark} & \textbf{Flink} \\
%                 \hline
%                 \textbf{Armazenamento} & Disco              & RAM            & RAM            \\
%                 \hline
%                 \textbf{Granularidade} & Grossa             & Grossa         & Fina           \\
%                 \hline
%                 \textbf{Estado}        & Sem                & Sem            & Com            \\
%                 \hline
%                 \textbf{Processamento} & Lote               & Micro lotes    & Stream         \\
%                 \hline
%                 \textbf{Volume}        & Finito             & Finito         & Infinito       \\
%                 \hline
%                 \textbf{Linguagem.}    & Java               & Scala          & Java           \\
%                 \hline
%             \end{tabular}
%         }
%         \caption*{\footnotesize Fonte: Elaborado pelo autor, 2023.}
%     \end{minipage}
% \end{table}

% \section{Citação}

% \subsection{Final do texto}
% \begin{verbatim}
% No atual cenário tecnológico, dados se tornaram um ativo de 
% alto valor \cite{gunther2017debating}.
% \end{verbatim}

% No atual cenário tecnológico, dados se tornaram um ativo de alto valor \cite{gunther2017debating}.

% \subsection{Início do texto}
% \begin{verbatim}
% Segundo \textcite{gunther2017debating} Dados se tornaram um 
% ativo de alto valor no atual cenário tecnológico.
% \end{verbatim}

% Segundo \textcite{gunther2017debating} Dados se tornaram um ativo de alto valor no atual cenário tecnológico.

% \section{Alíneas}
% Para criar alíneas utilize o comando enumerate, nunca description ou itemize. As alínea devem encerrar com um ponto e vírgula e a última deve encerrar com um ponto final.

% \begin{verbatim}
% \begin{enumerate}
%     \item Primeiro item da alínea;
%     \item Segundo item da alínea;
%     \item Terceiro item da alínea.
%     \begin{enumerate}
%         \item Primeiro item da subalínea;
%         \item Segundo item da subalínea;
%         \item Terceiro item da subalínea.
%     \end{enumerate}
% \end{enumerate}
% \end{verbatim}

% \begin{enumerate}
%     \item Primeiro item da alínea;
%     \item Segundo item da alínea;
%     \item Terceiro item da alínea.
%           \begin{enumerate}
%               \item Primeiro item da subalínea;
%               \item Segundo item da subalínea;
%               \item Terceiro item da subalínea.
%           \end{enumerate}
% \end{enumerate}

% \section{Customização}

% Esse pacote pode ser customizado passando argumentos da seguinte forma:
% \begin{verbatim}
% \usepackage[acronym, glossaries, index, labelref, debug]{CEFET}
% \end{verbatim}

% \begin{enumerate}
%     \item \textbf{acronym:} adiciona o suporte para lista de abreviaturas e siglas;
%     \item \textbf{glossaries:} adiciona o suporte para glossário;
%     \item \textbf{index:} adiciona o suporte para índice de assunto;
%     \item \textbf{labelref:} \verb|\ref{fig:1}| retorna Figura 1 em vez de 1 para todas as referências;
%     \item \textbf{debug:} Ativa as réguas e os quadros para melhorar a visualização das medidas.
% \end{enumerate}

% \section{Namedref}
% \begin{verbatim}
% Exemplo de utilização do \textbf{namedref} para a 
% \ref{tabela:ComparativoFrameworks}.
% \end{verbatim}
% Exemplo de utilização do \textbf{namedref} para a \ref{tabela:ComparativoFrameworks}.